% !PRINTED 128
\weberpage{128}
%<<CONT
the indices of $Q$ and of the divisors of $P$ conjugate to $Q$, taken with respect to $Q$. Hence $h_1 = 1$ and

\[
p^{\alpha-\beta} = h_1 + h_2 + h_3 \ldots % [sic: Weber prints no plus sign before the ellipsis; the preceding decomposition requires one.]
\]

The summands $h_1, h_2, h_3, \ldots$, being divisors of the degree of $Q$, can here only be powers of $p$; and since $h_1 = 1$; it follows that at least $p$ of these summands must be $= 1$, since otherwise their sum could not be divisible by $p$. But this means nothing else than that at least $p - 1$ of the conjugate groups $b^{-1}Qb$, where $b$ is not contained in $Q$, must be divisible by $Q$ and therefore equal to $Q$. But if

\[
b^{-1}Qb = Q
\]

then for every exponent $s$ we also have

\[
b^{-s}Qb^s = Q.
\]

If $b^r$ is the lowest power of $b$ contained in $Q$, then, since the degree of $b$ is a power of $p$, $r$ must also be a power of $p$, and

\[
b^{\frac{r}{p}} = c
\]

is an element of $P$ which does not itself occur in $Q$, but whose $p$th power is contained in $Q$, and for which at the same time the condition

\[
cQc^{-1} = Q
\]
holds.

From this it follows that the elements

\[
R = Q + Qc + Qc^2 + \cdots + Qc^{p-1}
\]

form a group of degree $p^{\beta+1}$, of which $Q$ is a normal divisor. We therefore have the theorem:

III. \gesperrt{If $P$ is a group of degree $p^\alpha$ and $Q$ a divisor of $P$ of degree $p^\beta$ ($\beta < \alpha$), then there is a divisor $R$ of $P$ of degree $p^{\beta+1}$, of which $Q$ is a normal divisor.}

If in this theorem one takes $\beta = \alpha - 1$, then $R$ coincides with $P$, and it results that $Q$ is a normal divisor of $P$. Applying the same theorem to the group $Q$, and so on, one obtains:

IV. \gesperrt{Every group whose degree is a power of a prime is metacyclic (\S.\,8).}
